% Source PDF: source/普通高考/2010/2010湖北文.pdf
% Page: 1; the source graph is vector content (pdfimages -list found no embedded bitmap).
% Elements: broken horizontal mass axis; vertical frequency/group-width axis; six contiguous histogram bars;
% dashed reference levels; x labels 0, 1.00--1.30 and y labels 0.4, 1, 3, 4, 5.6, 6.
% Relations: bins [1.00,1.05), [1.05,1.10), [1.10,1.15), [1.15,1.20),
% [1.20,1.25), [1.25,1.30) have frequency densities 1, 4, 5.6, 6, 3, 0.4.
% The compressed x coordinates below encode the printed axis break between 0 and 1.00.
\begin{tikzpicture}[baseline]
\begin{axis}[exam statistical histogram,
    width=15.053cm,
    height=10.161cm,
    axis lines=none,
    xmin=-0.9,
    xmax=8.4,
    ymin=-0.45,
    ymax=7.25,
    font=\normalsize,
    clip=false,
    ticks=none,
]
    % Dashed frequency/group-width reference levels, behind the bars.
    \draw[dashed,line width=0.55pt] (axis cs:0,0.4) -- (axis cs:7,0.4);
    \draw[dashed,line width=0.55pt] (axis cs:0,1) -- (axis cs:1,1);
    \draw[dashed,line width=0.55pt] (axis cs:0,3) -- (axis cs:5,3);
    \draw[dashed,line width=0.55pt] (axis cs:0,4) -- (axis cs:2,4);
    \draw[dashed,line width=0.55pt] (axis cs:0,5.6) -- (axis cs:4,5.6);
    \draw[dashed,line width=0.55pt] (axis cs:0,6) -- (axis cs:5,6);

    % Histogram bars, with unit-width compressed bins corresponding to 0.05 kg.
    \addplot[
        ybar interval,
        draw=black,
        fill=white,
        line width=0.75pt,
    ] coordinates {
        (1,1)
        (2,4)
        (3,5.6)
        (4,6)
        (5,3)
        (6,0.4)
        (7,0)
    };

    % Axes and the printed break on the mass axis.
    \draw[->,line width=0.8pt] (axis cs:0,0) -- (axis cs:8.05,0);
    \draw[->,line width=0.8pt] (axis cs:0,0) -- (axis cs:0,7.0);
    \examstatxbreak[line width=0.8pt]{0.6}{0};

    % x-axis ticks and labels.
    \draw[line width=0.45pt] (axis cs:0,0) -- (axis cs:0,-0.10);
    \draw[line width=0.45pt] (axis cs:1,0) -- (axis cs:1,-0.10);
    \draw[line width=0.45pt] (axis cs:2,0) -- (axis cs:2,-0.10);
    \draw[line width=0.45pt] (axis cs:3,0) -- (axis cs:3,-0.10);
    \draw[line width=0.45pt] (axis cs:4,0) -- (axis cs:4,-0.10);
    \draw[line width=0.45pt] (axis cs:5,0) -- (axis cs:5,-0.10);
    \draw[line width=0.45pt] (axis cs:6,0) -- (axis cs:6,-0.10);
    \draw[line width=0.45pt] (axis cs:7,0) -- (axis cs:7,-0.10);
    \node[anchor=north east,inner sep=0pt,font=\normalsize,yshift=-4pt] at (axis cs:-0.12,0.000000) {0};
    \node[anchor=north,inner sep=0pt,font=\normalsize,yshift=-2pt] at (axis cs:1,-0.100000) {1.00};
    \node[anchor=north,inner sep=0pt,font=\normalsize,yshift=-2pt] at (axis cs:2,-0.100000) {1.05};
    \node[anchor=north,inner sep=0pt,font=\normalsize,yshift=-2pt] at (axis cs:3,-0.100000) {1.10};
    \node[anchor=north,inner sep=0pt,font=\normalsize,yshift=-2pt] at (axis cs:4,-0.100000) {1.15};
    \node[anchor=north,inner sep=0pt,font=\normalsize,yshift=-2pt] at (axis cs:5,-0.100000) {1.20};
    \node[anchor=north,inner sep=0pt,font=\normalsize,yshift=-2pt] at (axis cs:6,-0.100000) {1.25};
    \node[anchor=north,inner sep=0pt,font=\normalsize,yshift=-2pt] at (axis cs:7,-0.100000) {1.30};

    % y-axis ticks and labels.
    \draw[line width=0.45pt] (axis cs:0,0.4) -- (axis cs:-0.16,0.4);
    \node[anchor=east,inner sep=0pt,font=\normalsize,xshift=-4pt] at (axis cs:-0.160000,0.4) {0.4};
    \draw[line width=0.45pt] (axis cs:0,1) -- (axis cs:-0.16,1);
    \node[anchor=east,inner sep=0pt,font=\normalsize,xshift=-4pt] at (axis cs:-0.160000,1) {1};
    \draw[line width=0.45pt] (axis cs:0,3) -- (axis cs:-0.16,3);
    \node[anchor=east,inner sep=0pt,font=\normalsize,xshift=-4pt] at (axis cs:-0.160000,3) {3};
    \draw[line width=0.45pt] (axis cs:0,4) -- (axis cs:-0.16,4);
    \node[anchor=east,inner sep=0pt,font=\normalsize,xshift=-4pt] at (axis cs:-0.160000,4) {4};
    \draw[line width=0.45pt] (axis cs:0,5.6) -- (axis cs:-0.16,5.6);
    \node[anchor=east,inner sep=0pt,font=\normalsize,xshift=-4pt] at (axis cs:-0.160000,5.6) {5.6};
    \draw[line width=0.45pt] (axis cs:0,6) -- (axis cs:-0.16,6);
    \node[anchor=east,inner sep=0pt,font=\normalsize,xshift=-4pt] at (axis cs:-0.160000,6) {6};

    \examstatylabel{\(\dfrac{\text{频率}}{\text{组距}}\)};
    \examstatxlabel{质量 (kg)};
\end{axis}
\end{tikzpicture}
